\section{Discussion}\label{sec:discussion}
The construction makes one part of synthetic population conditioning explicit: the rule translating aggregate preference signs into training labels. A policy fitted to those labels can score new item text without naming the population in the evaluation prompt. This is a useful computational object for studying how supplied group information affects synthetic responses. The restriction applies to the declared conditioning channel; the model's prior information and the semantic choices embedded in the shared bank remain part of the object.

The GPS--WVS application shows why construction success and external usefulness require separate evidence. Recovery of synthetic labels establishes that fitting moved the selected contrasts in the intended direction. Trust anchor ordering also transfers to a separate survey battery more strongly than in the country-named prompting comparator. Yet correspondence with human country ranks is uncertain. Proposed risk and patience indicators expose further discrepancies. We treat those discrepancies as questions about cross-survey coherence rather than as a diagnosis of which instrument failed. Convergent and discriminant evidence, together with cross-national comparability, would be needed to support a stronger construct-validity interpretation \cite{campbellfiske1959,davidov2014}. GPS country scores are the declared anchor. Human--GPS association is the criterion for already-fielded item maps. On the forty-two-country intersection, the nine-item trust map tracks GPS in the expected direction, with an interval that excludes zero. Adapter--GPS association is retention of that anchor on new wording. Adapter--human association is agreement with people and remains unresolved on the sixteen-country panel. Consistency with published development or institutional gradients can screen a candidate measure. Those screens do not replace the human columns.

The method could be applied to regions, cohorts, or other groups with comparable aggregate anchors. That is scope of the construction, while demonstrated empirical transport is limited to the selected countries and item bank. Sign-only labeling discards anchor intensity and makes countries with common sign profiles share a labeling problem. Joint fitting also prevents attributing a dimension-specific output uniquely to that dimension's coordinate. A magnitude-aware construction would be a different protocol; intensity weighting alone would not establish human predictive validity.

A prospective use is survey pre-piloting. Researchers could score candidate wording, compare its implied country ordering with declared anchors, and nominate items for independent human assessment. The present exercise does not establish that this process reduces cost, improves item selection, or shortens a questionnaire without losing validity. A blinded item-screening study should compare anchor-assisted selection with credible alternatives using a held-out human pilot. Appendix~\ref{app:humanscreen} records a prospective protocol for that study. Scoring one country or one item cannot certify measurement validity. Human observations also remain necessary for statistically justified rectification of synthetic estimates \cite{krsteski2026}.

Several limitations bound the evidence. The country panel is a census of two production waves rather than a probability sample of the GPS--WVS intersection, repeated sign profiles reduce distinct construction inputs, and a single shared bank does not test sensitivity to bank generation. The shared bank is mixed-trait. Intended loadings are generator judgments rather than experimentally verified trait manipulations, so off-target movement on non-target facets remains a construction risk for any later dimension-specific reading. We do not apply a trait-isolation filter, and we do not treat unused process diagnostics as evidence of purity. Prompt and checkpoint provenance differ across stored runs; bootstrap intervals condition on those runs rather than integrating their uncertainty. Neither pretraining exposure to survey material nor the accuracy of generator semantics is removed by country-token omission. The study also does not estimate effects of demographic or policy changes on human preferences. Perturbing an anchor tests the construction algorithm. Human causal claims require separate validation.

We intend the policies to complement field work. They should not be used to speak for individuals, replace confirmatory samples, or assign national cultural essences. A reusable scoring instrument can help formulate measurement questions, provided that its construction assumptions and unresolved validation tasks remain visible.
